\subsection{WEIGHTS AND PRIMITIVE ELEMENTS}

Let V be a g-module. For all \(\lambda\in h^{*}\) , denote by \(V^{\lambda}\) the primary subspace, relative to \(\lambda\) , of V considered as an h-module (Chap. VII, §1, no. 1). The elements of \(V^{\lambda}\) are called the elements of weight \(\lambda\) of the g-module V. The sum of the \(V^{\lambda}\) is direct (Chap. VII, §1, no. 1, Prop. 3). For all \(\alpha\in h^{*}\) and \(\lambda\in h^{*}\) , \(g^{\alpha}V^{\lambda}\subset V^{\alpha+\lambda}\) (Chap. VII, §1, no. 3, Prop. 10 (ii)). The dimension of \(V^{\lambda}\) is called the multiplicity of \(\lambda\) in V; if it is \(\geq1\) , i.e.~if \(V^{\lambda}\neq0\) , \(\lambda\) is said to be a weight of V. If V is finite dimensional, the homotheties of V defined by the elements of h are semi-simple, so \(V^{\lambda}\) is the set of \(x\in V\) such that \(Hx=\lambda(H)x\) for all \(H\in h\) .

Lemma 1. Let \(\mathrm{V}\) be a \(\mathfrak{g}\) -module and \(v \in \mathrm{V}\) . The following conditions are equivalent:

\begin{enumerate}
\def\labelenumi{(\roman{enumi})}
\item
  \(\mathfrak{b}_{+}v\subset kv;\)
\item
  \(\mathfrak{h}v \subset kv\) and \(n_{+}v = 0;\)
\item
  \(\mathfrak{h}v\subset kv\) and \(\mathfrak{g}^{\alpha}v = 0\) for all \(\alpha \in \mathrm{B}\) .
\item
  \(\Longrightarrow\) (ii): Assume that \(\mathfrak{b}_{+}v\subset kv\) . There exists \(\lambda \in \mathfrak{h}^*\) such that \(v\in V^{\lambda}\) . Let \(\alpha \in \mathrm{R}_{+}\) . Then \(\mathfrak{g}^{\alpha}.v\subset V^{\lambda}\cap V^{\lambda +\alpha} = 0\) . Hence \(\mathfrak{n}_{+}v = 0\) .
\item
  \(\Longrightarrow\) (iii): This is clear.
\item
  \(\Longrightarrow\) (i): This follows from the fact that \((X_{\alpha})_{\alpha \in \mathrm{B}}\) generates \(\mathfrak{n}_{+}\) (§3, no. 3, Prop. 9 (iii)).
\end{enumerate}

DEFINITION 1. Let V be a g-module and \(v \in V\) . Then v is said to be a primitive element of V if \(v \neq 0\) and v satisfies the conditions of Lemma 1.

A primitive element belongs to one of the \(V^{\lambda}\) . For all \(\lambda\in h^{*}\) , \(V_{\pi}^{\lambda}\) denotes the set of \(v\in V^{\lambda}\) such that \(b_{+}v\subset kv\) . Thus, the primitive elements of weight \(\lambda\) are the non-zero elements of \(V_{\pi}^{\lambda}\) .

PROPOSITION 1. Let V be a g-module, v a primitive element of V and \(\omega\) the weight of v. Assume that V is generated by v as a g-module.

\begin{enumerate}
\def\labelenumi{(\roman{enumi})}
\item
  If \(\mathrm{U}(\mathfrak{n}_{-})\) denotes the enveloping algebra of \(\mathfrak{n}_{-}\) , we have \(\mathrm{V} = \mathrm{U}(\mathfrak{n}_{-}).v\) .
\item
  For all \(\lambda \in \mathfrak{h}^*\) , \(V^{\lambda}\) is the set of \(x \in V\) such that \(Hx = \lambda(H)x\) for all \(H \in \mathfrak{h}\) . We have \(V = \bigoplus_{\lambda \in \mathfrak{h}^*} V^{\lambda}\) , and each \(V^{\lambda}\) is finite dimensional. The space \(V^{\omega}\) is of dimension 1, and every weight of \(V\) is of the form \(\omega - \sum_{\alpha \in B} n_{\alpha}. \alpha\) , where the \(n_{\alpha}\) are integers \(\geq 0\) .
\item
  V is an indecomposable \(\mathfrak{g}\) -module, and its commutant reduces to the scalars.
\item
  Let \(\mathrm{U}(\mathfrak{g})\) be the enveloping algebra of g, and Z the centre of \(\mathrm{U}(\mathfrak{g})\) . There exists a unique homomorphism \(\chi\) from Z to k such that, for all \(z \in Z\) , \(z_{V}\) is the homothety with ratio \(\chi(z)\) .
\end{enumerate}

Let \(\mathrm{U}(\mathfrak{b}_{+})\) be the enveloping algebra of \(\mathfrak{b}_{+}\) . We have \(\mathrm{U}(\mathfrak{g}) = \mathrm{U}(\mathfrak{n}_{-})\) . \(\mathrm{U}(\mathfrak{b}_{+})\) (Chap. I, §2, no. 7, Cor. 6 of Th. 1). Hence

\[
\mathrm{V} = \mathrm{U} (\mathfrak {g}). v = \mathrm{U} (\mathfrak {n} _ {-}). \mathrm{U} (\mathfrak {b} _ {+}). v = \mathrm{U} (\mathfrak {n} _ {-}). v.
\]

Denote by \(\alpha_{1},\ldots,\alpha_{n}\) the distinct elements of \(R_{+}\) . Then

\[
\left(X _ {- \alpha_ {1}} ^ {p _ {1}} X _ {- \alpha_ {2}} ^ {p _ {2}} \dots X _ {- \alpha_ {n}} ^ {p _ {n}}\right) (p _ {1}, \ldots , p _ {n}) \in \mathbf {N} ^ {n}
\]

is a basis of \(\mathrm{U}(\mathfrak{n}_{-})\) , so

\[
\mathrm{V} = \sum_ {(p _ {1}, \dots , p _ {n}) \in \mathbf {N} ^ {n}} k X _ {- \alpha_ {1}} ^ {p _ {1}} \dots X _ {- \alpha_ {n}} ^ {p _ {n}} v.\tag{1}
\]

For \(\lambda \in \mathfrak{h}^*\) , put

\[
\mathrm{T} _ {\lambda} = \sum_ {(p _ {1}, \ldots , p _ {n}) \in \mathbf {N} ^ {n}, \omega - p _ {1} \alpha_ {1} - \dots - p _ {n} \alpha_ {n} = \lambda} k X _ {- \alpha_ {1}} ^ {p _ {1}} \ldots X _ {- \alpha_ {n}} ^ {p _ {n}} v.
\]

By Chap. VII, §1, no. 1, Prop. 2 (ii), if \(h \in \mathfrak{h}\) , \(h_{\mathrm{V}}|_{\mathrm{T}_{\lambda}}\) is the homothety with ratio \(\lambda(h)\) . So \(\mathrm{T}_{\lambda} \subset \mathrm{V}^{\lambda}\) . On the other hand, (1) implies that

\[
\mathrm{V} = \sum_ {\lambda \in \omega - \mathbf {N} \alpha_ {1} - \dots - \mathbf {N} \alpha_ {n}} \mathrm{T} _ {\lambda}.
\]

The sum of the \(V^{\lambda}\) is direct (Chap. VII, §1, no. 1, Prop. 3). From these observations it follows that \(V^{\lambda} = T_{\lambda}\) , that V is the direct sum of the \(V^{\lambda}\) , and that \(V^{\lambda}\) is the set of \(x \in V\) such that \(hx = \lambda(h)x\) for all \(h \in h\) . On the other hand, \(\dim V^{\lambda}\) is at most the cardinal of the set of \((p_{1}, \ldots, p_{n}) \in \mathbf{N}^{n}\) such that \(p_{1}\alpha_{1} + \cdots + p_{n}\alpha_{n} = \omega - \lambda\) . This proves that \(V^{\lambda} = 0\) if \(\omega - \lambda \notin \sum_{\alpha \in B} N\alpha\) , that \(\dim V^{\omega} = 1\) , and that the \(V^{\lambda}\) are all finite dimensional.

Let \(c\) be an element of the commutant of V. For all \(h \in \mathfrak{h}\) ,

\[
h c (v) = c h (v) = \omega (h) c (v),
\]

so \(c(v) \in V^{\omega}\) ; hence there exists \(t \in k\) such that \(c(v) = tv\) . Now, for all \((p_1, \ldots, p_n) \in \mathbf{N}^n\) ,

\[
c X _ {- \alpha_ {1}} ^ {p _ {1}} \dots X _ {- \alpha_ {n}} ^ {p _ {n}} v = X _ {- \alpha_ {1}} ^ {p _ {1}} \dots X _ {- \alpha_ {n}} ^ {p _ {n}} c v = t X _ {- \alpha_ {1}} ^ {p _ {1}} \dots X _ {- \alpha_ {n}} ^ {p _ {n}} v
\]

so that \(c = t.1\) . Hence, the commutant of \(V\) reduces to the scalars. This implies (iv) and the fact that \(V\) is indecomposable.

DEFINITION 2. The homomorphism \(\chi\) of Prop. 1 (iv) is called the central character of the \(\mathfrak{g}\) -module V.

PROPOSITION 2. Let V be a g-module generated by a primitive element e of weight \(\omega\) , and X a semi-simple g-module. Let \(\Phi\) be the set of homomorphisms from the g-module V to the g-module X. Then \(\varphi\mapsto\varphi(e)\) is an isomorphism from the vector space \(\Phi\) to the vector space \(X_{\pi}^{\omega}\) .

It is clear that \(\varphi(e) \in X_{\pi}^{\omega}\) for all \(\varphi \in \Phi\) . If \(\varphi \in \Phi\) and \(\varphi(e) = 0\) , then \(\varphi = 0\) since \(e\) generates the \(\mathfrak{g}\) -module V. We show that, if \(f\) is a non-zero element of \(X_{\pi}^{\omega}\) , there exists \(\varphi \in \Phi\) such that \(\varphi(e) = f\) . Let \(X'\) be the submodule of X generated by \(f\) . By Prop. 1, \(X'\) is indecomposable, hence simple since X is semi-simple. The element \((e, f)\) is primitive in the \(\mathfrak{g}\) -module \(V \times X\) . Let N be the submodule of \(V \times X\) generated by \((e, f)\) . Then \(N \cap X \subset \mathrm{pr}_2(N) = X'\) , so \(N \cap X = 0\) or \(X'\) ; if \(N \cap X = X'\) , N contains the linearly independent elements \((e, f)\) and \((0, f)\) which are primitive of weight \(\omega\) ; this is absurd (Prop. 1), so \(N \cap X = 0\) . Thus \(\mathrm{pr}_1|_{N}\) is an injective map \(h\) from N to V; this map is surjective since its image contains \(e\) . Thus \(\varphi = \mathrm{pr}_2 \circ h^{-1}\) is a homomorphism from the \(\mathfrak{g}\) -module V to the \(\mathfrak{g}\) -module X such that \(\varphi(e) = f\) .

